\section{Modification recipes}
\label{sec:extending}

Each recipe names every file to touch, in order, and ends with the test
obligation. The general pattern is always the same: implement the physics
as a pure function in the right layer, thread it through the engine(s),
expose it in the GUI with a persisted variable, and pin it with a test.

\subsection{Add a filter}

\emph{Files touched: \code{data/filters.list}, \code{data/filters/}.}
For a filter that exists in the SVO Filter Profile Service, download the
VOTable of \emph{both} calibrations (the Vega page and the AB page of the
same filter), save them as \code{<System.Band>\_Vega.xml} and
\code{<System.Band>\_AB.xml} in \code{data/filters/}, and append the two
names to \code{data/filters.list}. No code changes: the selector is
rebuilt from the list at start-up and every number the program needs
(zero point, detector type, pivot, transmission) is inside the files.
For a filter SVO does not carry, produce the pair from a transmission
curve with \code{make\_filter\_profile.py} --- it prints the exact lines
to append. If the program refuses a profile, the parser's error message
names the missing \code{PARAM}; the four required ones are
\code{MagSys}, \code{ZeroPoint} (+\code{ZeroPointUnit} \code{Jy}),
\code{DetectorType} and \code{WavelengthUnit}.

\subsection{Add template spectra}

\emph{Files touched: \code{data/}, \code{data/interpola.db.csv} --- both
via the tool.} Run \code{python3 add\_template.py <file-or-directory>}.
The tool computes the one non-obvious catalogue field, $m_{v0}$ (the
synthetic Vega $V$ magnitude of the file --- get this wrong and every
prediction from that template is off by the same amount), writes the
row, and copies the file. Do it by hand only when you understand the
$m_{v0}$ convention (User Guide Sect.~6.1). To \emph{remove} a template,
delete its catalogue line; the spectrum file may stay.

\subsection{Add an observing site}

Use \textsf{Save current site} in the GUI, which writes
\code{etc\_sites.json}; or edit that JSON directly (schema in the User
Guide Sect.~6.7). \code{observatories.json} is the shipped default set:
add there only sites that should reach every user of a release.

\subsection{Add a noise term}

\emph{Files touched: \code{observing\_conditions.py},
\code{photometry.py}, \code{spectroscopy.py}, possibly
\code{solvers.py}, \code{etc\_gui.py}, \code{tests/}.} Write the term as
a pure function returning a \emph{variance in e$^{-2}$} given rates and
conditions, with the reference in the docstring. In
\code{photometry.compute\_photometry\_single} and in the S/N line of
\code{spectroscopy.compute\_spectroscopy}, add it to the
\code{extra\_variance} sums (grep for \code{scintillation\_var} to find
both spots). Ask the one question that decides solver integration: is
the variance linear in exposure time? If yes (like scintillation), also
add its rate to the \code{extra\_variance\_rate\_e2\_s} arguments of
\code{exposure\_time\_for\_snr} and \code{plan\_stack} at the call sites
in \code{etc\_gui.py} and \code{spetc\_batch.py}, or target-S/N mode
will silently ignore it. If it needs a new user input, register the Tk
variable with \code{self.\_var("key", default)} (persistence is then
automatic) and pass it through \code{\_atmosphere\_dict}. Finish with a
regression test pinned to a published magnitude of the effect and a line
in the physics document's noise table.

\subsection{Add a PSF model}

\emph{Files touched: \code{etc\_physics.py}, \code{etc\_gui.py},
\code{tests/}.} Add a branch to \code{psf\_encircled\_energy} (circular
integral) and \code{psf\_slit\_throughput} (the 1-D marginal CDF,
analytic if you can find it --- the Moffat case shows the pattern with
the Student-$t$ --- numeric otherwise; keep the \code{offset\_arcsec}
argument vectorised, the dispersion-loss code calls it with arrays).
Add the model name to the GUI combobox. Both engines pick it up
automatically because they read \code{psf\_model} from the atmosphere
dictionary. Test it against the Gaussian in a limiting case.

\subsection{Change or extend the sky model}

Component physics (a better zodiacal model, a site-specific airglow
coefficient) belongs in \code{sky\_brightness.py}; \emph{when} each
component sets the level belongs in \code{sky\_background.py}; how the
per-time dictionaries are assembled for the engines is
\code{\_sky\_models\_for\_track} in the GUI. A future fully spectral sky
model has a ready-made slot: put a curve under the
\code{spectral\_sky\_f\_lambda} key of the sky-model dictionary and the
spectroscopy engine uses it as-is, bypassing the broad-band colour
machinery.

\subsection{Add a GUI control}

Register the variable (\code{self.\_var}), place the widget in the right
column builder (\code{\_build\_*\_column}), consume the value in
\code{\_run\_etc} or a helper, and --- if the physics sees it --- extend
the run-parameters log in \code{\_show\_info} so the assumption is
visible in the output record. Never put computation in the GUI: compute
in an engine or module function and call it.

\subsection{Add an output quantity}

Compute it in the engine and put it in the result dict / DataFrame ---
it then reaches the result CSV automatically. To show it in the results
table, add a \code{(key, heading)} pair to
\code{ETCGUI.\_RESULT\_COLUMNS}; to show it per time sample, add a
column in \code{\_build\_time\_series\_dataframe}; to show it in the log,
add a line to the \code{\_append\_info\_outputs} call. Document it in
the README output table and the User Guide.
